\subsection{The automorphisms of the actual procongruence exit system}

\begin{result}{D.33}{Proved}{Genus zero classifies the procongruence exit system}
Let $\check{\mathcal T}$ be the ambient procongruence exit
system of D.31, with the same labelled strata, open objects,
boundary branches, actual twist embeddings, relabeling,
forgetting, nodal clutching and composition comparisons
as in E.42.  Equivalences are continuous and are taken
modulo coherent natural isomorphism.  Restriction gives
a topological group isomorphism
\[
       \operatorname{Aut}(\check{\mathcal T})
                       \xrightarrow{\sim}\widehat{\mathrm{GT}}.
                                                               \tag{D.48}
\]
The quotient from the full system to this procongruence
system induces the identity on GT under E.42 and (D.48).
This conclusion uses the actual cut-stabilizer extensions
in Boggi's Section 4 and the genus-at-most-two congruence
comparison, and uses no general profinite multitwist
normalizer assertion.  The additional outer group
branch-order comparison of E.37 selects exactly
$\operatorname{Fix}(\mathcal J)$ in both systems; E.44
proves this is the whole group through its explicit
ten-point forgetting construction.
\end{result}
\begin{proof}
We prove the full morphism comparisons needed to transfer
the kernel argument.  Throughout the proof a check over a
stabilizer means its closure in the ambient procongruence
group, not an unidentified full completion.

For an actual multicurve $\sigma$, the results used from
\cite{ET_Boggi2014} are precisely Proposition 3.3,
Proposition 3.4, Theorem 4.2, Theorem 4.5 and Remark 4.7,
all before Section 5.  They give the following exact maps.
Actual curve simplices embed in the procongruence curve
complex, and their stabilizers are the closures of their
discrete stabilizers.  For the subgroup fixing each curve
with its orientation, cutting gives
\[
 1\longrightarrow\widehat{\mathbb Z}^{\sigma}
 \longrightarrow\check C_{\vec\sigma}
 \longrightarrow\prod_v\check\Gamma_{g_v,n_v}
 \longrightarrow1.                                         \tag{D.49}
\]
The unpinned group is the actual extension
\[
 1\longrightarrow\check C_{\vec\sigma}
 \longrightarrow\check C_\sigma
 \longrightarrow A_\sigma\longrightarrow1,                 \tag{D.50}
\]
where $A_\sigma$ is the discrete stabilizer's finite image
on the oriented branches and components.  It is a subgroup
of the signed permutation group preserving the labelled
topological type.  These sequences are the actual closure,
cutting and permutation maps; no splitting is used.

For clarity, continuity and the exact finite image in
(D.50) can be checked on the finite set of oriented curve
classes $[\gamma^{\pm1}]$ in the Hausdorff space of profinite
conjugacy classes.  Their setwise stabilizer acts
continuously on this finite set.  Density gives precisely
the discrete image $A_\sigma$, and its kernel is the
closure of the discrete finite-index pinned subgroup.
Likewise, for several retained subsets of a fixed actual
multicurve, their joint stabilizer is the inverse image
of the corresponding stabilizer of this finite flag.
Thus its closure is exactly that joint stabilizer.
This computes the closures of actual flag groups,
including their endpoint projections.

The link comparison in Remark 4.7 is the join of the
procongruence curve complexes of the cut components,
with action given by (D.49).  One can also see its
precise completion step directly.  A profinite simplex
having $\sigma$ as a face lies in one of finitely many
ambient simplex orbits.  Choose an actual representative
and an actual transport of its specified face to
$\sigma$.  The remaining transport lies in
$\check C_\sigma$, by the preceding stabilizer equality,
and is a limit of actual transports fixing $\sigma$.
Hence the simplex is a limit of actual simplices in
that link.  Cutting these is the discrete join
comparison.  The surjection onto the full product in
(D.49) supplies exactly the product congruence topology,
proving the stated join and its action.  In particular
normalizing a retained actual multicurve neither loses
nor adds an incidence orbit.

We now compute every completed Hom-space from the finite
ambient quotients defining $\check{\mathcal T}$.  Write
$\Gamma=\Gamma_{g,n}$ and choose its actual skeletal
multicurves $\sigma_i$.  For every submulticurve
$\rho\subseteq\sigma_i$ in the orbit of $\sigma_j$,
choose $g_\rho$ with $g_\rho\rho=\sigma_j$.  The
discrete transporter is the finite disjoint union
$\coprod_\rho C_{\sigma_j}g_\rho$, as in (E.177).
Each $C_{\sigma_j}$ is congruence-separable: an actual
mapping class in its closure stabilizes the embedded
actual simplex, and therefore lies in the discrete
stabilizer.  Thus any two distinct displayed cosets
are separated by a finite congruence quotient.
Since right source twists preserve each coset,
taking the inverse limit gives
\[
 \begin{aligned}
 \check\Hom(i,j)&=
 \coprod_\rho
       \check C_{\sigma_j}g_\rho/\widehat\Delta_{\sigma_i}\\
 &\lhook\joinrel\longrightarrow
       \check\Gamma/\widehat\Delta_{\sigma_i}.
 \end{aligned}                                               \tag{D.51}
\]
Here the right quotient in a summand uses its conjugate
twist subgroup inside $\check C_{\sigma_j}$.  All these
lattices are embedded by (D.49).  The map on the right
is the actual composition with $[1]:j\to\varnothing$.
It is injective: the closure of each transporter
coset has the indicated embedded quotient, and a
cofinal set of the finite congruence levels separates
the finitely many distinct cosets.  Thus no two
summands or elements are identified in the limit.  The endpoint group is
$\check C_{\sigma_j}/\widehat\Delta_{\sigma_j}$.
The multiplication of representatives supplies all
composition maps in (D.51).

These formulas give a direct congruence version of the
framed-square extraction in E.40.  Attach the same
freshly labelled tripods and let $t$ be their tag
multicurve.  The automorphism group of its full
crossing is exactly $\check C_t$, by (D.51), with
its target inclusion and source quotient.  This
defines the framed procongruence core group
$\check\Gamma^{\partial}$, including its full
boundary twist lattice.  Its quotient by that
lattice is the actual capped group, by (D.49).
Mapping classes between different marked core
objects give the corresponding closure torsors;
restriction and extension by the tripods are their
inverse comparison maps.

Allow arbitrary interior core multicurves
$\gamma_0\supseteq\gamma_1$.  As in (E.160), use
the two vertical crossings
$t\cup\gamma_i\to\gamma_i$.
The new label pairs distinguish $t$ from all
essential interior core curves.  Thus the two
intermediate subsets in the two factorizations
of their common composite are intrinsic.
After normalizing at $\gamma_1$, (D.49)--(D.51)
identify the total-composite orbit and both its
factorization orbits.  Their stabilizers are the
same actual flag closure, computed by (D.50);
there are no additional intersections of unknown
full-completion subgroups.  The finite-orbit
compactness argument of E.40 applies directly:
the finitely many compact endpoint orbits contain
the dense product of discrete Hom-spaces and
therefore cover that whole product.  Equations
(E.166)--(E.168), with the actual check groups
in place of the full groups, then give the
square space
\[
       (\check C_1\times\check C_0)/\check S,
 \quad C_i=\Stab(t,\gamma_i)/\Delta_{\gamma_i},\quad
 S=\Stab(t,\gamma_0,\gamma_1)/\Delta_{\gamma_1}.              \tag{D.52}
\]
The check groups in this display are the closures
and quotients furnished by normalization and
(D.49).  The link comparison above identifies
their core incidences with the component
procongruence incidences.  Hence (D.52) is
exactly the completed framed-core transporter,
with both horizontal projections.  Multiplication
agrees on the dense discrete products, so it
also gives its full composition comparison.
We have re-proved the square equalizer in the
congruence category; we have not obtained it by
quotienting an equalizer of full completions.

It remains to verify smooth sewing on these
whole Hom-spaces.  Use E.42's three-node chain
$i,m,j$ with its two distinct companion labels.
Equation (E.183) shows that the two labelled
$(0,3)$ vertices cannot be created by a
nonempty contraction.  Thus $i,m,j$ are
individually intrinsic, and the square proof
just given applies to these chain flags as
well.  Retain every remaining outer tag.

Let $\check L_i$ be its diagonal group at a
core object, and $\check P_i$ the framed core
group there.  Cutting at the middle node is
an actual map $q:\check L_i\to\check P_i$.
After this normalization, the companion and
the new $m$-branch are the two fresh labels
of each original tripod tag; the retained
$i,j$ and core flags therefore give precisely
the framed groups of (D.49)--(D.52).
By (D.49) its kernel is exactly
$\widehat{\mathbb Z}_m$, and its image is all
of $\check P_i$: on the discrete groups it
is the projection $P_i\times\mathbb Z_m\to P_i$
of (E.184), so compactness and density give
this image assertion.  The kernel is central:
$m$ is individually fixed and all the surface
maps preserve orientation, hence conjugation
fixes its positive twist.  Normalize instead at
$i,j$ and select the middle four-point
component.  Its crossing group is
$\widehat{\mathbb Z}_m$, since genus zero
has full congruence topology.  This gives
$r:\check L_i\to\widehat{\mathbb Z}_m$
which is the identity on the kernel of $q$.
Consequently
\[
             (q,r):\check L_i\xrightarrow{\sim}
                    \check P_i\times\widehat{\mathbb Z}_m.
                                                               \tag{D.53}
\]
Injectivity follows from the stated kernel;
surjectivity follows by lifting a $q$-coordinate
and correcting its $r$-coordinate by the
middle twist.  This proves the splitting by
the two actual maps, without inserting a
section into the cutting extension.

The identical argument applies to every
incidence stabilizer $\check S_L$: its
$q$-image is $\check S_P$, its kernel is the
same middle twist, and $r$ restricts to the
identity there.  Hence
$\check S_L=\check S_P\times\widehat{\mathbb Z}_m$
under (D.53).  The homogeneous square spaces
are therefore exactly
\[
 \begin{aligned}
 &\frac{(\check P_1\times\widehat{\mathbb Z})
             \times(\check P_0\times\widehat{\mathbb Z})}
       {\check S_P\times\operatorname{diag}(\widehat{\mathbb Z})}\\
 &\qquad\simeq
       ((\check P_1\times\check P_0)/\check S_P)
                          \times\widehat{\mathbb Z}.
 \end{aligned}                                               \tag{D.54}
\]
by $[(p_1,n_1),(p_0,n_0)]\mapsto([p_1,p_0],n_1-n_0)$.
Its inverse sets $n_0=0$.  This proves $q,r$ on
every framed Hom-space, all endpoint actions
and all compositions.  Actual finite branch
permutations are retained through (D.50).

The simultaneous zero-$r$ section is now
intrinsic and unique on Hom-spaces, in the
isomorphism-closed fibre formulation of
(E.187).  Forget the two companions and
stabilize.  The actual formula
$q_e=q_iq_mq_j$ of (E.188) sends that section
to smooth sewing with zero middle displacement.
For several sewings there is one $r$ for each
middle node; the unique simultaneous section
proves the same associativity and self-sewing
coherence.  Marked caps and unmarked disc caps
are the actual forgetting operations of D.31,
with C.1's stabilization maps.  The annular
unit is the one-node four-point crossing.
Thus the complete whole-arrow sewing
construction of E.42 is recovered in the
procongruence system, with every comparison
map now proved for its actual topology.

We finish with the kernel.  Restrict an
equivalence of $\check{\mathcal T}$ to the
recovered marked-pants sewing groupoids.
Their genus-zero part is the same as its
full completion by D.19.  The strict
object-fixing kernel proof in E.27 applies
without a change of any group or word:
every group used in that proof is of
genus zero or one, and its framed version
is a tripod stabilizer in the same genus.
D.19, together with (D.49), identifies
all those check groups, torsors, injections
and sewing maps with the full groups
used there.  In particular its actual
$\widehat B_3\to\widehat\Gamma_{1,2}$
injection, its two cut subgroups, and
the (G2) computation are unchanged.
E.27 therefore fixes the one-holed-torus
move $\sigma$ for a genus-zero kernel
element.  The other three generating
moves are in genus zero.  Their discrete
operadic composites generate all marked
mapping classes, and are dense in the
defined procongruence torsors.  Continuity
forces the strict automorphism to be
the identity in every genus.  Coherent
object strictification uses the same
free marked-pants object operad and
cyclic relabelings as E.35.

For each nodal endpoint, normalize its
actual labelled graph, use this coherent
identity comparison on its cut-component
groupoids, and apply nodal clutching.
Relabeling coherence descends it through
the actual finite groupoid (D.50).
This fixes one comparison for the endpoint,
independently of its incident arrows.
Equations (D.53)--(D.54) and the unique
zero section have these same endpoint
comparisons, since $q,r$ are precisely
the stated cutting and component maps.

In those coordinates the four-point
crossing is fixed.  The explicit smooth
self-sewing and disc cap
$(0,4)\to(1,2)\to(1,1)$ in (E.189)
fix the one-holed-torus crossing too.
Cutting any adjacent incidence at its
retained curves, taking the two pants
adjacent to its smoothed curve, and
sewing back all other cuffs expresses
that actual pinned crossing by these
operations, as proved in E.42.  A
change of pants refinement is a sequence
of already fixed component morphisms,
so it keeps the same endpoint comparisons.
Every adjacent crossing and every
endpoint translate are consequently fixed.
Every completed arrow is a composite
of adjacent ones: there are finitely
many intermediate object lists, and
the compact images of their composition
maps contain the dense discrete Hom-space.
The normalized equivalence is therefore
the identity on the entire category,
with its operation comparisons.  This
proves injectivity in (D.48).

D.31 supplies an inverse by descending
the full equivalence $\mathcal F_A$ of
E.39.  Both restriction and this inverse
are continuous on the finite quotient
systems, by their actual proword and
transporter formulas.  They are the
same restriction and GT parameter as
in E.42, proving the asserted comparison
of the two automorphism groups.  Finally
the additional outer branch-order
diagram of E.37 consists entirely of
the low-genus groups just identified;
its computed condition remains exactly
$A=\mathcal J(A)$.  No conclusion about
vanishing of an ambient higher-genus
congruence kernel is part of this
automorphism comparison; those kernels
are the groups explicitly measured in
D.32.
\end{proof}
